\documentclass[12pt, twoside]{article}
\input{/Users/chris/github/course-files/docs/ib/preamble}

\fancyhead[LE]{\thepage}
\fancyhead[RO]{Markscheme}
\fancyhead[LO]{La Scuola d'Italia / Dr.\ Huson / IB Math 12 HL \\* 17 November 2026}

\begin{document}
\subsubsection*{3.6 Quiz: Periodic Functions (HL) --- Markscheme}

\emph{Total: 61 marks. Problems 1--3 are without a calculator: exact answers
are required. In Problems 4--6, accuracy marks need 3 s.f., the exact value,
or the calculator display with an ellipsis; a value cut short with no
ellipsis earns A0 (markscheme-spec \S21). Problems 4--6 need radian mode: a
value worked in degree mode earns the method marks only. Problems 5 and 6
are answered on lined paper.}

\begin{enumerate}[itemsep=0.3cm]

%--- Problem 1: f(x) = -4cos(3x) - 1; p sin form; g = f(x + pi/4) [10 marks] ---
\item[1.] [Maximum mark: 10]

$f(x) = a\cos(bx) + d$; minimum $(0, -5)$, maximum
$\left(\frac{\pi}{3}, 3\right)$; $g$ is $f$ translated $\frac{\pi}{4}$ to
the left.

\textbf{(a)} The graph starts at a minimum, so $a < 0$:

$a = -\dfrac{3 - (-5)}{2} = -4$ \hfill \textbf{A1}

$d = \dfrac{3 + (-5)}{2} = -1$ \hfill \textbf{A1}

\emph{Wrong method: $a = 4$ (the sign not read from the graph): A0.}

\textbf{(b)} minimum to maximum is half a period, so the period is
$2 \times \dfrac{\pi}{3} = \dfrac{2\pi}{3}$ \hfill \textbf{M1}

$b = \dfrac{2\pi}{2\pi/3} = 3$ \hfill \textbf{A1}

\emph{Wrong methods: $b = 6$ ($\frac{\pi}{3}$ taken as the period):
M0A0; $b = \frac{2}{3}$ (period used as $b$): M1A0.}

\textbf{(c)} $p = 4$ \hfill \textbf{A1}

$-4\cos 3x = 4\sin\left(3x - \dfrac{\pi}{2}\right)
= 4\sin\left(3\left(x - \dfrac{\pi}{6}\right)\right)$ \hfill \textbf{M1}

$q = \dfrac{\pi}{6}$ \hfill \textbf{A1}

\emph{M1 for any valid method: $\sin\left(\theta - \frac{\pi}{2}\right) =
-\cos\theta$; or the sine reaching its maximum where
$3(x - q) = \frac{\pi}{2}$ at $x = \frac{\pi}{3}$; or the sine crossing the
principal axis upward a quarter period before the maximum. Wrong method:
$q = \frac{\pi}{2}$ (from $3x - \frac{\pi}{2}$ without factoring the $3$):
M1A0. FT from (a) and (b).}

\newpage

\textbf{(d)(i)} $g(x) = f\!\left(x + \dfrac{\pi}{4}\right)
= -4\cos\left(3\left(x + \dfrac{\pi}{4}\right)\right) - 1$ \hfill \textbf{A1}

\emph{Accept $-4\cos\left(3x + \frac{3\pi}{4}\right) - 1$ and
$4\sin\left(3\left(x + \frac{\pi}{12}\right)\right) - 1$. Wrong methods,
A0: $f\!\left(x - \frac{\pi}{4}\right)$ (the graph moved left, not right);
$-4\cos\left(3x + \frac{\pi}{4}\right) - 1$.}

\textbf{(d)(ii)} $g(0) = -4\cos\dfrac{3\pi}{4} - 1$ \hfill \textbf{M1}

$= -4\left(-\dfrac{\sqrt{2}}{2}\right) - 1 = 2\sqrt{2} - 1$ \hfill \textbf{A1}

\emph{M1 for substituting $x = 0$ into the candidate's $g$. The answer must
be exact: $1.83$ earns A0. FT from (d)(i).}

\newpage

%--- Problem 2: sec(theta) = -4 in quadrant III; cos 105 deg; sec(7pi/12) [10 marks] ---
\item[2.] [Maximum mark: 10]

$\sec\theta = -4$, $\pi < \theta < \frac{3\pi}{2}$ (third quadrant).

\textbf{(a)} $1 + \tan^{2}\theta = \sec^{2}\theta = 16$ \hfill \textbf{M1}

$\tan^{2}\theta = 15$ \hfill \textbf{A1}

$\tan\theta = \sqrt{15}$ \quad (tangent is positive in the third quadrant)
\hfill \textbf{A1}

\emph{Accept the method $\cos\theta = -\frac{1}{4}$, $\sin\theta =
-\frac{\sqrt{15}}{4}$, then the quotient. The last A1 needs the positive
value only: $-\sqrt{15}$ or $\pm\sqrt{15}$ earns M1A1A0. A bare $\sqrt{15}$
earns M0A1A1.}

\textbf{(b)} $\cos\theta = -\dfrac{1}{4}$, \quad
$\sin\theta = -\sqrt{1 - \dfrac{1}{16}} = -\dfrac{\sqrt{15}}{4}$
\hfill \textbf{M1}

$\operatorname{cosec}\theta = -\dfrac{4}{\sqrt{15}}$ \hfill \textbf{A1}

\emph{Accept $-\frac{4\sqrt{15}}{15}$. Accept $1 + \cot^{2}\theta =
\operatorname{cosec}^{2}\theta$ with $\cot\theta = \frac{1}{\sqrt{15}}$ from
(a). The A1 needs the negative sign (sine is negative in the third
quadrant). FT from (a).}

\textbf{(c)} $\cos 105^\circ = \cos 60^\circ\cos 45^\circ - \sin 60^\circ\sin 45^\circ$
\hfill \textbf{M1}

$= \dfrac{1}{2}\cdot\dfrac{\sqrt{2}}{2} - \dfrac{\sqrt{3}}{2}\cdot\dfrac{\sqrt{2}}{2}
= \dfrac{\sqrt{2} - \sqrt{6}}{4}$ \hfill \textbf{A1}

\emph{Accept $\frac{1 - \sqrt{3}}{2\sqrt{2}}$. Wrong methods: a plus sign in
the compound-angle formula gives $\frac{\sqrt{2} + \sqrt{6}}{4}$ (that is
$\cos 15^\circ$): M0A0; $\cos 60^\circ + \cos 45^\circ$: M0A0.}

\textbf{(d)} $105^\circ = 105 \times \dfrac{\pi}{180} = \dfrac{7\pi}{12}$
\hfill \textbf{A1}

$\sec\dfrac{7\pi}{12} = \dfrac{4}{\sqrt{2} - \sqrt{6}}
= \dfrac{4(\sqrt{2} + \sqrt{6})}{(\sqrt{2} - \sqrt{6})(\sqrt{2} + \sqrt{6})}$
\hfill \textbf{M1}

$= \dfrac{4(\sqrt{2} + \sqrt{6})}{2 - 6} = -\sqrt{2} - \sqrt{6}$
\hfill \textbf{A1}

\emph{M1 for multiplying by the conjugate. Accept $-(\sqrt{2} + \sqrt{6})$.
Wrong method: $\sec = \frac{1}{\sin}$: M0A0. FT from (c).}

\newpage

%--- Problem 3: cos 2x - 3 sin x - 2 = 0 on [0, 2pi] [7 marks] ---
\item[3.] [Maximum mark: 7]

\textbf{(a)} $1 - 2\sin^{2}x - 3\sin x - 2 = 0$ \hfill \textbf{M1}

$-2\sin^{2}x - 3\sin x - 1 = 0$ \hfill \textbf{A1}

$2\sin^{2}x + 3\sin x + 1 = 0$ \hfill \textbf{AG}

\emph{M1 for substituting $\cos 2x = 1 - 2\sin^{2}x$. Using
$2\cos^{2}x - 1$ or $\cos^{2}x - \sin^{2}x$ earns the M1 only once it is
converted to sine. Wrong method: $\cos 2x = 1 - 2\sin x$: M0A0.}

\textbf{(b)} $(2\sin x + 1)(\sin x + 1) = 0$ \hfill \textbf{M1}

$\sin x = -\dfrac{1}{2}$ \ or \ $\sin x = -1$ \hfill \textbf{A1}

$x = \dfrac{7\pi}{6}$, \ $x = \dfrac{11\pi}{6}$, \ $x = \dfrac{3\pi}{2}$
\hfill \textbf{A1A1A1}

\emph{M1 for factorising or the quadratic formula in $\sin x$. A1 for each
value of $x$. If extra values are given (e.g.\ $\frac{\pi}{6}$,
$\frac{5\pi}{6}$ from $\sin x = \frac{1}{2}$, or $-\frac{\pi}{6}$ outside
the interval), deduct one A mark. Answers in degrees ($210^\circ$,
$330^\circ$, $270^\circ$) earn M1A1 and A0 for the values. Wrong method:
dividing through and losing $\sin x = -1$ loses the A1 for
$\frac{3\pi}{2}$. FT from the candidate's values of $\sin x$.}

\newpage

%--- Problem 4: Ambiguous case: AB = 10, AC = 7, angle B = 0.6 rad [9 marks] ---
\item[4.] [Maximum mark: 9]

$AB = 10$, $AC = 7$, $A\hat{B}C = 0.6$ radians, $BC = x$.

\textbf{(a)} $7^{2} = 10^{2} + x^{2} - 2(10)(x)\cos 0.6$ \hfill \textbf{M1}

$49 = 100 + x^{2} - 20x\cos 0.6$ \hfill \textbf{A1}

$x^{2} - (20\cos 0.6)\,x + 51 = 0$ \hfill \textbf{AG}

\emph{M1 for the cosine rule with $AC = 7$ opposite the angle $0.6$. A1 for a
correct line before the given result.}

\textbf{(b)} $x = 12.4$ \quad ($12.3907\ldots$) \quad and \quad
$x = 4.12$ \quad ($4.11597\ldots$) \hfill \textbf{A1A1}

\emph{GDC polynomial solver or the quadratic formula. Wrong method: degree
mode gives $17.0$ and $3.00$: A0A0.}

\textbf{(c)} $\dfrac{\sin A\hat{C}B}{10} = \dfrac{\sin 0.6}{7}$ \hfill \textbf{M1}

$A\hat{C}B = 0.938$ \quad ($0.938431\ldots$) \hfill \textbf{A1}

$A\hat{C}B = \pi - 0.938431\ldots = 2.20$ \quad ($2.20316\ldots$)
\hfill \textbf{A1}

\emph{Accept the cosine rule with each value of $x$. Both angles are
possible: $0.6 + 2.20316\ldots < \pi$. Answers in degrees ($53.8^\circ$,
$126^\circ$) earn M1A0A0. Wrong method: giving only $0.938$ earns M1A1A0.}

\textbf{(d)} $\frac{1}{2}(10)(12.3907\ldots)\sin 0.6 - \frac{1}{2}(10)(4.11597\ldots)\sin 0.6$
\hfill \textbf{M1}

$= 34.9816\ldots - 11.6202\ldots = 23.4\,\text{cm}^{2}$ \quad ($23.3613\ldots$)
\hfill \textbf{A1}

\emph{M1 for two area expressions subtracted (any correct form, e.g.\ using
$\frac{1}{2}(7)(x)\sin A\hat{C}B$). Accept any value consistent with the
candidate's correctly rounded $x$ values, e.g.\ $23.4$ from $12.4$ and
$4.12$. Wrong method: one triangle's area only: M0A0. FT from (b).}

\newpage

%--- Problem 5: f(x) = sin x + sqrt3 cos x = 2 sin(x + pi/3) [12 marks] ---
\item[5.] [Maximum mark: 12]

\textbf{(a)} $2\sin\left(x + \dfrac{\pi}{3}\right)
= 2\left(\sin x\cos\dfrac{\pi}{3} + \cos x\sin\dfrac{\pi}{3}\right)$
\hfill \textbf{M1}

$= 2\left(\dfrac{1}{2}\sin x + \dfrac{\sqrt{3}}{2}\cos x\right)$
\hfill \textbf{A1}

$= \sin x + \sqrt{3}\cos x$ \hfill \textbf{AG}

\textbf{(b)} a translation of $\dfrac{\pi}{3}$ to the left (horizontal
translation by $-\frac{\pi}{3}$) \hfill \textbf{A1}

a vertical stretch with scale factor $2$ \hfill \textbf{A1}

\emph{Either order. Words are required (markscheme-spec \S19): each
transformation needs its type, direction and size. Accept ``translate by
$\binom{-\pi/3}{0}$''. Reject ``move''. Wrong method: translation to the
right: A0 for that transformation.}

\textbf{(c)} $-2 \le f(x) \le 2$ \hfill \textbf{A1}

\emph{Accept $[-2,\ 2]$. Strict inequalities earn A0.}

\textbf{(d)} $2\sin\left(x + \dfrac{\pi}{3}\right) = 1$, so
$\sin\left(x + \dfrac{\pi}{3}\right) = \dfrac{1}{2}$ \hfill \textbf{M1}

$\dfrac{\pi}{3} \le x + \dfrac{\pi}{3} \le \dfrac{7\pi}{3}$, so
$x + \dfrac{\pi}{3} = \dfrac{5\pi}{6}$ or $\dfrac{13\pi}{6}$ \hfill \textbf{A1}

$x = \dfrac{\pi}{2}$ \hfill \textbf{A1}

$x = \dfrac{11\pi}{6}$ \hfill \textbf{A1}

\emph{The first A1 for at least one correct value of $x + \frac{\pi}{3}$ in
the shifted interval. Wrong method: using $x + \frac{\pi}{3} = \frac{\pi}{6}$
gives $x = -\frac{\pi}{6}$, outside the domain; if it is given as an answer,
deduct one A mark. GDC decimals ($1.57$, $5.76$) earn M1A1A0A0.}

\textbf{(e)} $x + \dfrac{\pi}{3} = \pi$ or $2\pi$, so $x = \dfrac{2\pi}{3}$
and $x = \dfrac{5\pi}{3}$ \hfill \textbf{A1}

$\text{Area} = \left|\displaystyle\int_{2\pi/3}^{5\pi/3}
2\sin\left(x + \dfrac{\pi}{3}\right)\mathrm{d}x\right|
= \left|\Bigl[-2\cos\left(x + \tfrac{\pi}{3}\right)\Bigr]_{2\pi/3}^{5\pi/3}\right|$
\hfill \textbf{M1}

$= \left|-2\cos 2\pi + 2\cos\pi\right| = 4$ \hfill \textbf{A1}

\emph{Accept $2.09$ and $5.24$ for the intercepts. M1 for an integral of $f$
(or $|f|$) between the candidate's intercepts; a GDC integral is accepted.
Accept $-4$ as the integral with ``area $= 4$''. FT from the candidate's
intercepts.}

\newpage

%--- Problem 6: Daylight in New York and Palermo; rate of change; days P > N [13 marks] ---
\item[6.] [Maximum mark: 13]

$N(t) = a\sin\bigl(b(t - c)\bigr) + d$; maximum $15.1$ at $t = 172$,
minimum $9.3$ at $t = 355$; Palermo maximum $14.8$, minimum $9.5$.

\textbf{(a)} $a = \dfrac{15.1 - 9.3}{2} = 2.9$, \quad
$d = \dfrac{15.1 + 9.3}{2} = 12.2$ \hfill \textbf{A1A1}

\emph{Bare values earn A1A1.}

\textbf{(b)} $b = \dfrac{2\pi}{365}$ \quad ($0.0172142\ldots$) \hfill \textbf{A1}

$\dfrac{2\pi}{365}(172 - c) = \dfrac{\pi}{2}$ \hfill \textbf{M1}

$c = 172 - 91.25 = 80.75$ \hfill \textbf{A1}

\emph{Accept $c = 81.25$ from the minimum. Wrong method: $c = 172$ (the
cosine shift): M0A0.}

\textbf{(c)} $p = 2.65$, \quad $q = 12.15$ \hfill \textbf{A1A1}

\textbf{(d)} $N'(t) = 2.9 \times \dfrac{2\pi}{365}\cos\left(\dfrac{2\pi}{365}(t - 80.75)\right)$
\hfill \textbf{M1}

$N'(321) = -0.0272$ hours per day \quad ($-0.0272190\ldots$) \hfill \textbf{A1}

$-0.0272190\ldots \times 60 = -1.63$ minutes per day \quad ($-1.63314\ldots$)
\hfill \textbf{A1}

\emph{M1 for differentiating with the chain rule, or for a GDC derivative at
$t = 321$. Accept ``decreasing at $1.63$ minutes per day''. Accept any value
consistent with the candidate's rounded $b$ and $c$ (e.g.\ $-1.65$ from
$c = 81.25$; $-1.64$ from $b = 0.0172$). Wrong method: degree mode gives
$+2.99$: M1A0A0. FT from (a) and (b).}

\textbf{(e)} $P(t) = N(t)$ \hfill \textbf{M1}

\textit{GDC: graph $y = P(t) - N(t)$, or both models, for $0 \le t \le 365$.}

$t = 69.1$ \ ($69.0528\ldots$) \ and \ $t = 275$ \ ($274.947\ldots$)
\hfill \textbf{A1}

$P(t) > N(t)$ before $t = 69.0528\ldots$ and after $t = 274.947\ldots$:

$365 - (274.947\ldots - 69.0528\ldots) = 159$ days \quad ($159.105\ldots$)
\hfill \textbf{A1}

\emph{($P - N = -0.25\sin\bigl(b(t - c)\bigr) - 0.05$: Palermo has more
daylight from early October to early March.) Accept $159$ from
$b = 0.0172$ ($158.935\ldots$). Wrong method: $206$ days (the time when New
York has more daylight): M1A1A0. FT from the candidate's models.}

\end{enumerate}
\end{document}
